\UnnumberedChapter{Epilogue: genomic computation in three senses}{Epilogue}
\label{book:epilogue}\setcounter{figure}{0}\renewcommand{\thefigure}{E.\arabic{figure}}

The preface separated three meanings of genomic computation. One is computation
\emph{about} DNA, where recorded sequence is data. Another is computation \emph{with}
DNA, where molecules are the machine. The third is computation \emph{inspired by} DNA,
where genomic organization is a design hypothesis for digital systems. This book has
worked almost entirely in the first and third senses. This epilogue steps back. It asks
what the frontier in each sense has actually demonstrated, how the book's evidence
discipline applies to each, and where the three might meet. It treats these meeting
points as research hypotheses, not as forecasts.

\BackPlate{three-meanings}{Three senses of genomic computation. Each has its own
objects, its own notion of a successful computation and its own strongest form of
evidence. The labelled interfaces between columns (sequencing and synthesis, design
hypotheses, typed backends) are where the senses could meet, and each needs its own
contract and test before it carries a claim.}{book:threemeanings}

\section*{Computation about DNA: what the frontier models teach}
Sequence models of genomes have changed quickly. DNABERT adapted masked pretraining to
overlapping $k$-mer tokens~\cite{dnabert2021}, and DNABERT-2 made tokenization efficiency
itself an architectural question~\cite{dnabert2}. Enformer predicted gene-expression
and chromatin tracks from sequence by integrating information over distances of up to
roughly a hundred kilobases~\cite{enformer2021}. Nucleotide Transformer showed the value
of scale, data diversity and a systematic downstream evaluation
protocol~\cite{nt2025}. HyenaDNA asked whether a non-attention operator could reach
long single-nucleotide contexts~\cite{hyenadna}. Caduceus built reverse-complement
equivariance and bidirectionality into the architecture~\cite{caduceus}. Evo and
Evo~2 pushed single-nucleotide modelling to genome scale and across all domains of
life~\cite{evo2024,evo2}. AlphaGenome predicts thousands of functional-genomic tracks,
up to single-base resolution, from a one-megabase input window, and in its authors'
evaluation it matched or exceeded the strongest external models on most of their
variant-effect benchmarks~\cite{alphagenome2026}.

Several of this book's chapters exist because these models made their questions
unavoidable. Tokenization changes what a prediction means (Chapters~5 and~24). Strand
symmetry is an architectural choice with a testable transformation law, not a
data-augmentation habit (Chapters~18 and~21). Long context is first an access graph
and only then a speed claim (Chapter~24). The hardest lessons concern evidence.
Benchmark comparisons between DNA foundation models turn out to be task-dependent, and
frozen representations and fine-tuned competitors answer different
questions~\cite{evod48:benchmark}. Homologous and overlapping sequences can leak across
a random split (Chapters~25 and~48). A model's sensitivity to a mutation is not the
mechanism a cell uses (Chapters~17 and~48).

Against this frontier, DOGMA and Hermon DNA are small and honest. They are reference
implementations with exact contracts, not trained foundation models, and their pilots
stayed at or near chance on a synthetic motif task under the declared budgets. The scientific
value of the two directions is that their claims are now stated precisely enough to
test. The three experiments that would most change the assessment are these:
\begin{enumerate}
\setlength{\itemsep}{1pt}
\item a matched-budget comparison of DOGMA's addressable loci with a capacity-matched
recurrent or state-space baseline on tasks with sparse, delayed dependencies;
\item a Hermon DNA ablation that holds the Transformer backbone fixed and varies only a
genomic mechanism (span-aware positions, tied strand features or local-plus-global
access) on homology-aware splits;
\item a study of the Evolutor runtime on persistent workflows that measures recovery,
provenance and policy compliance together with task success.
\end{enumerate}
Each comes with a falsifier written before the run, as Chapters~22, 29 and~51 require.

\begin{EvolutorNote}{A boundary this book keeps}
Some frontier genomic models are also used to \emph{generate} biological sequences. This
book does not cover designing sequences intended to function in living systems.
Its examples are synthetic, its evaluation designs concern prediction, and its
biological follow-up protocols stop at stating hypotheses and controls. Generative
biological design raises biosafety and governance questions that need specialist
oversight. Those questions fall outside an implementation textbook.
\end{EvolutorNote}

\section*{Computation with DNA: what molecules can now compute}
In 1994 Adleman solved a small instance of a directed Hamiltonian path problem with DNA
strands, ligation and selection~\cite{adleman1994}. The field that followed now has a
programmable primitive in toehold-mediated strand displacement, whose kinetics can be
tuned over many orders of magnitude~\cite{zhangwinfree2009}. It also has a compiler
precedent: any formal chemical reaction network can, in principle, be approximated by
DNA strand-displacement reactions~\cite{soloveichik2010}. Strand-displacement cascades
compute digital circuits~\cite{qianwinfree2011} and implement neural
networks~\cite{qianwinfreebruck2011}. Molecular winner-take-all networks recognize
patterns of many bits~\cite{cherryqian2018}. Algorithmic self-assembly runs
reprogrammable molecular algorithms from one tile set~\cite{woods2019}, and the
nucleation kinetics of self-assembly can itself classify patterns~\cite{evans2024}.
DNA also stores digital data at high density, with error-correcting codes designed for
its synthesis and sequencing channel~\cite{church2012,goldman2013,erlich2017}.

Two recent results changed what ``DNA computing'' can mean. In 2025, Cherry and Qian
reported supervised learning carried out by DNA itself. Training examples were written
into molecular concentration memories that then served as the weights for classifying
test patterns, all in vitro~\cite{cherryqian2025}. In the same year, Song and Qian showed
heat-rechargeable circuits and neural networks that compute repeatedly with new
inputs~\cite{songqian2025}. Learning and reuse, long the weakest points of molecular
computation, now have concrete demonstrations.

The resource semantics of Chapter~16 apply directly, and they expose how different the
substrate is. A molecular computation runs on very many molecules at once, but each
bimolecular step at nanomolar concentration takes minutes to hours. Its ``memory'' is a
set of concentrations with leak and decay, and its output is a fluorescence or
sequencing measurement with noise, not a register read. A fair comparison with
electronic computation must state units, workload and error model, exactly as
Chapters~29 and~47 demand of silicon. On those terms molecular computers are not general
replacements for electronic ones. Their distinctive strengths are operating where
electronics cannot (inside a test tube or a sample), interfacing directly with
biomolecules, and storing information densely and durably.

\section*{Where the senses could meet}
Evolutor's typed contracts suggest several concrete interfaces between the three
senses. Each is stated as a hypothesis with the evidence it would need.

\begin{description}
\item[DNA as an archival memory class.] Chapter~40 separated native state, durable
events, lossy summaries and addressable records. DNA storage could be a further class:
very dense and long-lived, but slow to write and read and subject to a characteristic
error channel. The hypothesis is that some evidence capsules and lineage records
(Chapters~43 and~52) are written once and read rarely, so a durable molecular tier
could hold them. Evidence would need a typed contract (capacity, write and read latency,
retrieval granularity, error and erasure semantics) and a cost comparison with
conventional archives under the same retention requirement. Deletion is especially
hard, which matters for the forgetting semantics of Chapter~40.
\item[A molecular backend behind the runtime.] Chapters~39 and~41 route typed operations
to backends whose contracts are declared. A strand-displacement or self-assembly
program could in principle be such a backend for operations that must happen in a
sample. Its contract would include reaction semantics, concentration units,
time-to-result, leak and error rates, single-use versus rechargeable
resources~\cite{songqian2025} and the readout's measurement model. The compiler
precedent already exists, from formal reaction networks to DNA
reactions~\cite{soloveichik2010}. The open questions are verification and reliability
at scale, together with an honest account of when a wet backend is worth its latency.
\item[Models that read molecular programs.] Sequence models could help \emph{design and
check} molecular computers: predicting hybridization, secondary structure and leak
for strand-displacement components before synthesis. This keeps the model in the
first sense (computation about DNA) and the machine in the second. It requires
validation against measured kinetics, not against the model's own predictions.
\item[Molecular learning as a comparison point.] Supervised learning in DNA neural
networks~\cite{cherryqian2025} gives Part~VI's methodology an unusual test case. The
same estimands, splits, controls and uncertainty apply, and so do the same warnings
about selecting on the test set, whether the learner is silicon or chemistry.
\item[Inspiration that must earn its keep.] The third sense, the book's own DOGMA and
Evolutor directions, borrows organization rather than chemistry. The relevant evidence
is never that a mechanism resembles biology. It is a matched-budget comparison in which
removing the borrowed organization costs something measurable.
\end{description}

\section*{What would count as evidence of general or superhuman capability}
Some readers will hope that systems built from these ideas lead towards artificial
general intelligence, or towards capabilities beyond human performance across broad
domains. Chapter~51 explained why no finite test establishes unrestricted generality,
and why breadth, depth and autonomy form a vector rather than a
level~\cite{evod51:levels}. The same discipline gives five questions to ask of any claim
of general or superhuman capability, including claims made about this book's
directions:
\begin{enumerate}
\setlength{\itemsep}{1pt}
\item \textbf{Which task families, and how were they held out?} Name the population of
tasks, show that development never saw the final suite, and audit pretraining exposure.
\item \textbf{Against which human or system baseline, under what budget?}
``Superhuman'' needs a defined comparison group, matched information and time, and
reported uncertainty.
\item \textbf{How efficiently is new skill acquired?} Report quality as a function of
data, compute, actions and assistance, not only final scores (Chapter~51).
\item \textbf{What happens out of distribution and under adversarial pressure?}
Include the failure cases, deferral behaviour and calibration, not only the successes.
\item \textbf{Is it reproducible and safe to operate?} Independent replication, policy
compliance, isolation, revocation and human control need their own evidence. They do
not follow from capability scores (Chapters~43--45).
\end{enumerate}
None of this book's results answers these questions. That is its scope, not a verdict
on the questions. A research program that wants to answer them will need trained
models at scale, sealed evaluations and independent replication. It will also need the
habits this book has tried to make automatic: name the mechanism, derive the
consequence, and test the claim in a way that could have failed.

\section*{Open problems}
The problems below collect the most consequential unresolved questions raised across
the book. Each is phrased so that a result could settle it in either direction.
\begin{enumerate}
\setlength{\itemsep}{1pt}
\item \textbf{Addressable versus compressed memory.} For a fixed task-information
budget, when do explicit loci or retained KV beat a compressed state, measured as
quality per state byte and per unit latency?
\item \textbf{Regulation beyond gating.} Does any regulation-inspired mechanism improve
quality or efficiency beyond ordinary input-dependent gating and routing after capacity
matching?
\item \textbf{Strand symmetry in practice.} On which genomic tasks does exact
reverse-complement equivariance improve sample efficiency, and on which does it erase
orientation-specific signal?
\item \textbf{Long-context genomics.} Which distal dependencies do learned models
actually use? Answer with distal-context ablations on homology-aware splits, not with
attention maps.
\item \textbf{Exact engines at scale.} Can the exact-refinement proofs of Parts~VII
and~VIII (state pools, paged KV, speculative decoding) be carried into device kernels
with tested completion fences and measured benefit?
\item \textbf{Effects in compilers.} What effect system lets an AI-workflow compiler
reorder and cache tool, memory and model calls safely, and how is that verified?
\item \textbf{Continual adaptation without regression.} Can structural and parametric
updates be promoted under the governance of Chapter~43 with measured, bounded
regressions?
\item \textbf{Persistent agents.} Do provenance-aware memory and intent-before-dispatch
recovery reduce human intervention on long-horizon tasks without increasing policy
violations?
\item \textbf{Molecular tiers.} Under what retention and access patterns does a
DNA-based archival tier beat conventional archives on total cost and durability?
\item \textbf{Learning across substrates.} Can molecular learners and electronic learners
be compared on a shared task with shared estimands and controls, and what does the
comparison reveal about each substrate's limits?
\end{enumerate}
The book ends where research begins: with precise objects, honest negative results and
experiments that have not yet been run.
